\section{Related Work}
\label{sec:related_work}

To the author's knowledge, no existing tool converts Markdown into an H5P Interactive Book.
Three groups of existing work are related to it.
General document converters could have been extended instead of developing a separate converter.
Plain-text question formats that Moodle already imports offer an alternative to the question extension.
Finally, other H5P content types could have served as the target instead of Interactive Book.

\paragraph{Document Converters.}
Pandoc~\cite{Pandoc} is a general-purpose document converter.
It reads several Markdown dialects and writes dozens of output formats.
A new output format for Pandoc is a writer that works on its document model, so an H5P writer could
have built on a mature parser and a complete model of the input.

For PDF input, a converter into H5P already exists.
\code{pdf\_to\_h5p\_converter} renders every page of a PDF as an image and turns each image into a
slide of a Course Presentation~\cite{Anandh:PdfToH5P}.
Rendering to images would have been possible here as well and much simpler, since only images are
emitted and the constraints on HTML and block placement do not arise.
However, text in an image cannot be searched, selected, resized, or read by a screen reader. It cannot
contain links or questions, and it cannot be edited in the H5P editor, which rules out the editing cycle supported by the merge feature.
The two tools thus represent opposite choices: preserving the appearance of the source or
preserving its structure.

\paragraph{Plain-Text Question Formats.}
If only questions are needed, Moodle can already import them from plain text.
The GIFT format~\cite{Moodle:GIFT} describes a question and its answers in a line-based syntax with
curly braces, and the Aiken format~\cite{Moodle:Aiken} is a stricter format for multiple-choice questions
with lettered answers.
In either case, the questions are imported into the question bank through Moodle's import page~\cite{Moodle:ImportQuestions}.
The converter's own question syntax follows a similar approach but is simpler than either format.

\paragraph{Other H5P Content Types.}
Besides Interactive Book, other H5P content types could have served as the target.
Course Presentation divides content into slides, but a slide is a fixed-size area on which every
element is positioned absolutely.
Text longer than a slide has to be split by hand, so a chapter of
lecture notes does not fit this model.
Page, by contrast, is a single column of content corresponding to one Markdown file, but a book of several files would
become several unrelated pages without navigation between them.
Question Set is a sequence of questions without room for text.
Interactive Book matches the structure of a folder of Markdown files: each file becomes a chapter,
content flows down a column instead of being positioned, and the reader gets a table of contents and
a progress indicator.
Its chapters also accept the layout and interactive content types needed for pseudo-tables and the
editing cycle.
Choosing Interactive Book determined most of the converter's constraints, since the
accepted content types and their \code{semantics.json} define what a Text block may contain.
